% ============================================================
% 05-bezier-matrix-sympy-text-output-explained.tex
% Compile with:  pdflatex -shell-escape 05-bezier-matrix-sympy-text-output-explained.tex
% Run twice for table of contents.
% ============================================================

\documentclass[12pt, a4paper]{article}

% ---- Encoding & fonts ----------------------------------------
\usepackage[T1]{fontenc}
\usepackage[utf8]{inputenc}
\usepackage{lmodern}

% ---- Page geometry -------------------------------------------
\usepackage[
  top=2.5cm, bottom=2.5cm,
  left=2.8cm, right=2.8cm
]{geometry}

% ---- Colours -------------------------------------------------
\usepackage[dvipsnames]{xcolor}
\definecolor{codebg}{RGB}{248,248,248}
\definecolor{rulecolor}{RGB}{180,180,180}
\definecolor{examplebg}{RGB}{240,248,255}
\definecolor{examplerule}{RGB}{100,149,237}
\definecolor{notebg}{RGB}{255,250,230}
\definecolor{noterule}{RGB}{210,180,50}
\definecolor{diffbg}{RGB}{240,255,240}
\definecolor{diffrule}{RGB}{80,160,80}

% ---- Minted --------------------------------------------------
\usepackage{minted}
\setminted[text]{
  bgcolor    = codebg,
  frame      = leftline,
  framerule  = 1.5pt,
  rulecolor  = rulecolor,
  fontsize   = \small,
  breaklines = true,
}

% ---- Math ----------------------------------------------------
\usepackage{amsmath}
\usepackage{amssymb}

% ---- Boxes ---------------------------------------------------
\usepackage{mdframed}

\newmdenv[
  backgroundcolor  = examplebg,
  linecolor        = examplerule, linewidth=1.5pt,
  leftline=true, rightline=false, topline=false, bottomline=false,
  innerleftmargin=10pt, innerrightmargin=8pt,
  innertopmargin=6pt, innerbottommargin=6pt,
  skipabove=6pt, skipbelow=4pt,
]{examplebox}

\newmdenv[
  backgroundcolor  = notebg,
  linecolor        = noterule, linewidth=1.5pt,
  leftline=true, rightline=false, topline=false, bottomline=false,
  innerleftmargin=10pt, innerrightmargin=8pt,
  innertopmargin=6pt, innerbottommargin=6pt,
  skipabove=6pt, skipbelow=4pt,
]{notebox}

\newmdenv[
  backgroundcolor  = diffbg,
  linecolor        = diffrule, linewidth=1.5pt,
  leftline=true, rightline=false, topline=false, bottomline=false,
  innerleftmargin=10pt, innerrightmargin=8pt,
  innertopmargin=6pt, innerbottommargin=6pt,
  skipabove=6pt, skipbelow=4pt,
]{diffbox}

% ---- Hyperlinks ----------------------------------------------
\usepackage[colorlinks=true, linkcolor=black, urlcolor=MidnightBlue]{hyperref}

% ---- Misc ----------------------------------------------------
\usepackage{booktabs}
\usepackage{needspace}
\usepackage{parskip}
\usepackage{titlesec}
\usepackage{fancyhdr}
\usepackage{datetime2}

% ---- Section formatting --------------------------------------
\titleformat{\section}
  {\large\bfseries\color{MidnightBlue}}
  {\thesection.}{0.6em}{}[\vspace{-4pt}\rule{\linewidth}{0.4pt}]
\titleformat{\subsection}
  {\normalsize\bfseries\color{BrickRed}}
  {\thesubsection.}{0.5em}{}

% ---- Header / footer -----------------------------------------
\pagestyle{fancy}
\fancyhf{}
\rhead{\textcolor{gray}{\small 05-bezier-matrix-sympy-text-output-explained}}
\lhead{\textcolor{gray}{\small B\'ezier Matrix Symbolic --- Output Explanation}}
\cfoot{\thepage}
\renewcommand{\headrulewidth}{0.3pt}

% ---- Helpers -------------------------------------------------
\newcommand{\cd}[1]{\texttt{#1}}

% ---- Title ---------------------------------------------------
\title{%
  \vspace{1cm}
  {\LARGE\bfseries B\'ezier Curve --- Matrix Method (Symbolic)}\\[0.4em]
  {\large $P(u) = U \cdot M \cdot G$ \quad via Symbolics.jl}\\[0.4em]
  {\large Explanation of the Program Output}\\[0.4em]
  {\normalsize Case 2: Quartic 3-D, 5 Control Points}\\[0.4em]
  {\normalsize\texttt{05-bezier-matrix-sympy-text.jl}}
}
\author{E.R.\ Nurwijayadi}
\date{\DTMtoday}


% =============================================================
\begin{document}
% =============================================================

\maketitle
\thispagestyle{empty}

\begin{center}
\textit{Directed and curated by E.R.\ Nurwijayadi.
Code and explanations AI-generated.}
\end{center}

\begin{center}
\begin{minipage}{0.85\textwidth}
\small
This document explains the output of
\cd{05-bezier-matrix-sympy-text.jl} when run with \cd{CASE = 2}.
This script is a hybrid: it uses the straight output structure
of script~04 (five flat sections, no derivation steps)
combined with exact \cd{Rational\{Int\}} arithmetic throughout,
and Symbolics.jl for polynomial formatting.
The output is almost identical to script~04 with three
specific differences.
\end{minipage}
\end{center}

\vspace{1em}\hrule\vspace{1.5em}
\tableofcontents
\newpage


% =============================================================
\needspace{8\baselineskip}
\section{What Changed from Script 04}
% =============================================================

\textit{Script~05 is script~04 with exact arithmetic substituted
for floating-point.
The output structure is identical --- same five sections,
same table format.
The differences are surgical.}

\begin{diffbox}
\textbf{Three differences from script~04:}

\medskip
\begin{tabular}{lp{4.5cm}p{4.5cm}}
\toprule
Location & Script~04 & Script~05 \\
\midrule
$u$ values in table &
  Decimals: \cd{0.2500}, \cd{0.5000}, \cd{0.7500} &
  Exact fractions: \cd{1//4}, \cd{1//2}, \cd{3//4} \\[6pt]
Resulting Polynomials section &
  Plain strings: \cd{13u\^{}4 - 32u\^{}3 + 24u\^{}2} &
  Symbolics.jl format: \cd{(24//1)*(u\^{}2) - (32//1)*(u\^{}3) + (13//1)*(u\^{}4)} \\[6pt]
Section subtitle &
  \cd{Resulting Polynomials} &
  \cd{Resulting Polynomials (via Symbolics.jl)} \\
\bottomrule
\end{tabular}

\medskip
The $G$, $M$, and $M \cdot G$ matrices look identical
because all their entries happen to be whole numbers.
\cd{Rational\{Int\}} displays integers without the \cd{//1}
suffix when the denominator is 1.
\end{diffbox}


% =============================================================
\needspace{8\baselineskip}
\section{Header}
% =============================================================

\begin{minted}{text}
==================================================================
  Bezier Curve - Matrix Method  [P(u) = U * M * G]  (Symbolics.jl)
  Case 2  |  Quartic  |  3D  |  5 Control Points
==================================================================
\end{minted}

The title now includes \cd{(Symbolics.jl)} to flag
that the polynomial formatting uses the CAS.
Everything else is unchanged from script~04.


% =============================================================
\needspace{8\baselineskip}
\section{$G$, $M$, and $M \cdot G$ --- Unchanged in Appearance}
% =============================================================

\begin{minted}{text}
  G  - Geometry Matrix (Control Points)
      x  y  z  ...  P0  0  0  1  ...  P4  5  4  1

  M  - Basis Matrix (derived via binomial)
       P0   P1   P2  P3  P4
  u^4   1   -4    6  -4   1  ...

  M * G  - Polynomial Coefficient Matrix
         x    y  z
  u^4   13  -28  0  ...
\end{minted}

These three sections look byte-for-byte identical to script~04.
The reason is that all values in $G$, $M$, and $M \cdot G$
are whole numbers.
\cd{Rational\{Int\}} stores them as fractions internally
(e.g.\ \cd{13//1}), but the formatting function \cd{fmt\_rat}
checks whether the denominator is 1 and prints only the
numerator in that case --- so \cd{13//1} displays as \cd{13},
not \cd{13//1}.

\begin{notebox}
\textbf{The subtitle change:} \cd{M --- Basis Matrix (derived via binomial)}

Script~04 said \cd{derived from Bernstein}.
Script~05 says \cd{derived via binomial}.
Both mean the same thing --- the entries come from the
binomial theorem applied to the Bernstein expansion ---
but script~05 is slightly more precise about the mechanism.
\end{notebox}


% =============================================================
\needspace{8\baselineskip}
\section{Resulting Polynomials --- Symbolic Format}
% =============================================================

\begin{minted}{text}
  Resulting Polynomials  (via Symbolics.jl)
  ----------------------------------------------------------------

  x(u)  =  (24//1)*(u^2) - (32//1)*(u^3) + (13//1)*(u^4)
  y(u)  =  (16//1)*u - (48//1)*(u^2) + (64//1)*(u^3) - (28//1)*(u^4)
  z(u)  =  1
\end{minted}

This section shows the most visible change from script~04.
The polynomial strings are now produced by Symbolics.jl
rather than by the hand-written \cd{fmt\_poly} function.

\needspace{5\baselineskip}
\subsection{Reading the Symbolics.jl notation}

The notation \cd{(24//1)*(u\^{}2)} has been explained in the
script~02 output explanation.
To recap:

\begin{itemize}
  \item \cd{24//1} is the \cd{Rational\{Int\}} value
        $\frac{24}{1} = 24$, printed with the rational type
        made explicit.
  \item \cd{*(u\^{}2)} is Symbolics.jl's way of writing
        multiplication by $u^2$.
  \item So \cd{(24//1)*(u\^{}2)} means exactly $24u^2$.
\end{itemize}

\begin{examplebox}
\textbf{Script~04 vs Script~05 for $x(u)$:}

\medskip
Script~04 (hand-formatted):
\begin{minted}{text}
  x(u)  =  13u^4  - 32u^3  + 24u^2
\end{minted}

Script~05 (Symbolics.jl formatted):
\begin{minted}{text}
  x(u)  =  (24//1)*(u^2) - (32//1)*(u^3) + (13//1)*(u^4)
\end{minted}

Two things differ:
\begin{enumerate}
  \item \textbf{Coefficient notation:} plain \cd{13} vs
        \cd{(13//1)}.
        Both mean the same number; the rational wrapper
        is Symbolics.jl being explicit about the type.
  \item \textbf{Term ordering:} script~04 prints from
        highest power down ($u^4$ first).
        Symbolics.jl prints from lowest power up
        ($u^2$ first) --- this is its default canonical ordering.
\end{enumerate}

The polynomials are mathematically identical.
\end{examplebox}

\begin{notebox}
\textbf{Why does $z(u)$ still print as just \cd{1}?}

The $z$ column of $M \cdot G$ is $[0, 0, 0, 0, 1]$ ---
the constant polynomial.
When Symbolics.jl builds the symbolic expression, it gets
$0 \cdot u^4 + 0 \cdot u^3 + \cdots + 1 \cdot u^0 = 1$.
A constant has no variable, so there is no \cd{(1//1)*u\^{}0}
--- Symbolics.jl simplifies it to the plain integer \cd{1}.
\end{notebox}


% =============================================================
\needspace{8\baselineskip}
\section{Evaluation Table --- Exact $u$ Values}
% =============================================================

\begin{minted}{text}
  Evaluation Table  -  P(u) = U * (M * G)
  ----------------------------------------------------------------

  +=============================================================+
  | u       x(u)               y(u)               z(u)        |
  |=============================================================|
  | 0       0  (0.0000)        0  (0.0000)   1  (1.0000)      |
  | 1//4    269//256  (1.0508) 121//64 (1.8906) 1  (1.0000)   |
  | 1//2    45//16  (2.8125)   9//4  (2.2500)   1  (1.0000)   |
  | 3//4    1053//256 (4.1133) 201//64 (3.1406) 1  (1.0000)   |
  | 1       5  (5.0000)        4  (4.0000)       1  (1.0000)  |
  +=============================================================+
\end{minted}

\needspace{5\baselineskip}
\subsection{Exact $u$ values}

Script~04 showed \cd{0.2500}, \cd{0.5000}, \cd{0.7500}
in the $u$ column.
Script~05 shows \cd{1//4}, \cd{1//2}, \cd{3//4}.

This is because \cd{u\_values} is now a
\cd{Vector\{Rational\{Int\}\}}, not a \cd{Vector\{Float64\}}.
The values $\tfrac{1}{4}$, $\tfrac{1}{2}$, $\tfrac{3}{4}$
are stored as exact fractions, and \cd{fmt\_rat} displays them
as \cd{1//4}, \cd{1//2}, \cd{3//4}.

\needspace{5\baselineskip}
\subsection{Exact result values}

The result values (\cd{269//256}, \cd{45//16}, etc.) are
the same fractions as in script~02.
The key difference from script~04 is \emph{how} they are
computed:

\begin{center}
\begin{tabular}{lp{5cm}p{5cm}}
\toprule
Script & How exact fractions are produced \\
\midrule
Script~04 &
  Compute with \cd{Float64}, then call \cd{rationalize()} to
  convert back to the nearest fraction.
  The fraction is recovered, not computed. \\[6pt]
Script~05 &
  Compute with \cd{Rational\{Int\}} throughout.
  The fraction \cd{269//256} is the direct result of
  integer arithmetic --- never a float. \\
\bottomrule
\end{tabular}
\end{center}

Both scripts display the same fractions.
Script~05 arrives at them by a more rigorous route.

\begin{examplebox}
\textbf{Verify $x\!\left(\tfrac{1}{4}\right) = \dfrac{269}{256}$
using \cd{Rational\{Int\}} arithmetic:}

The $U$ vector at $u = \tfrac{1}{4}$ is:
\[
  U = \left[\left(\tfrac{1}{4}\right)^4,\;
            \left(\tfrac{1}{4}\right)^3,\;
            \left(\tfrac{1}{4}\right)^2,\;
            \tfrac{1}{4},\; 1\right]
  = \left[\tfrac{1}{256},\; \tfrac{1}{64},\;
          \tfrac{1}{16},\; \tfrac{1}{4},\; 1\right]
\]

The $x$ column of $M \cdot G$ is $[13, -32, 24, 0, 0]$.

Dot product (all as exact fractions):
\[
  \tfrac{1}{256}(13) + \tfrac{1}{64}(-32) + \tfrac{1}{16}(24)
  + \tfrac{1}{4}(0) + 1(0)
\]
\[
  = \tfrac{13}{256} - \tfrac{128}{256} + \tfrac{384}{256}
  = \tfrac{13 - 128 + 384}{256}
  = \tfrac{269}{256}
\]

No floating-point involved at any step. \checkmark
\end{examplebox}

\needspace{5\baselineskip}
\subsection{Decimal column}

The decimal column (\cd{1.0508}, \cd{2.8125}, etc.)
is produced by converting the exact \cd{Rational\{Int\}}
result to \cd{Float64} at the very last moment:
\cd{Float64(269//256) = 1.05078125}, displayed as \cd{1.0508}
with 4 decimal places.
The decimal is a presentation convenience, not the actual
stored value.


% =============================================================
\needspace{8\baselineskip}
\section{The Convergence of All Five Scripts}
% =============================================================

\textit{Five scripts. Three methods. Two arithmetic systems.
One curve.}

\begin{center}
\begin{tabular}{llll}
\toprule
Script & Method & Arithmetic & Polynomial display \\
\midrule
01 & Bernstein expansion & \cd{Float64} & Hand-formatted \\
02 & Bernstein expansion & \cd{Rational\{Int\}} & Symbolics.jl \\
03 & Matrix (pedagogical) & \cd{Float64} & Hand-formatted \\
04 & Matrix (straight) & \cd{Float64} & Hand-formatted \\
05 & Matrix (straight) & \cd{Rational\{Int\}} & Symbolics.jl \\
\bottomrule
\end{tabular}
\end{center}

Every script produces the same five curve points:
$(0,0,1)$, $(1.051, 1.891, 1)$, $(2.813, 2.250, 1)$,
$(4.113, 3.141, 1)$, $(5, 4, 1)$.

The method (Bernstein vs matrix), the arithmetic
(\cd{Float64} vs \cd{Rational\{Int\}}), and the display
format (hand-formatted vs Symbolics.jl) are independent
choices that can be mixed freely.
Script~05 completes the matrix side of the
\cd{Rational\{Int\}} + Symbolics.jl combination,
mirroring what script~02 did for the Bernstein side.

\vspace{2em}\hrule\vspace{0.5em}
{\small\textcolor{gray}{%
Output explanation for \texttt{05-bezier-matrix-sympy-text.jl}, Case~2
\textbar{} Directed and curated by E.R.\ Nurwijayadi
\textbar{} Code and explanations AI-generated}}

\end{document}
